% Fichier complété par tools/complete_missing_corrections.py.
% Source : STAT/STAT-03-Demarche/08_RR3D/08_RR3D.tex
% Statut : rédigé (3 question(s)).
\begin{corrigebox}[Corrigé rédigé]
\CorrigeQuestion{1}
Le graphe relie 0--1 puis 1--2 par les deux pivots. On ajoute les
                poids, les couples des actionneurs et les réactions de liaison; les actions
                internes apparaissent par paires opposées.
\CorrigeQuestion{2}
Chaque pivot idéal transmet trois composantes de force et deux
                composantes de moment, sauf celle suivant son axe. Le poids de chaque
                solide est $\{-m_i g\vec j_0;\vec0\}_{G_i}$ et l'actionneur applique un
                couple pur suivant l'axe du pivot.
\CorrigeQuestion{3}
On isole 2 et on écrit le moment statique sur l'axe 2/1, puis on isole
                1+2 et on écrit le moment sur l'axe 1/0. Ces deux équations éliminent les
                réactions et donnent les deux couples moteurs requis.
\end{corrigebox}
